\chapter{Menstrual Disorders}

\section*{Definition and Overview}
Menstrual disorders encompass a broad range of conditions affecting the normal menstrual cycle, which is typically characterised by a cycle length of 24 to 38 days, a duration of 4 to 8 days of bleeding, and a normal blood loss of less than 80 mL per cycle. These disorders include heavy menstrual bleeding (menorrhagia), painful menstruation (dysmenorrhoea), absent menstruation (amenorrhoea), infrequent menstruation (oligomenorrhoea), and irregular intermenstrual bleeding. Menstrual abnormalities can arise from structural uterine abnormalities, hormonal imbalances, systemic coagulopathies, or lifestyle factors. Because they affect reproductive-aged individuals during their most productive years, menstrual disorders significantly impact physical health, quality of life, and emotional wellbeing, necessitating accurate diagnosis, individualised treatment, and fertility considerations.

\section*{Epidemiology}
Menstrual disorders are highly prevalent, affecting up to 75\% to 90\% of reproductive-aged individuals at some point. Dysmenorrhoea is the most frequently reported symptom, with a prevalence of 70\% to 90\% in adolescent girls and young women, and is a leading cause of short-term school and work absenteeism. Heavy menstrual bleeding is also extremely common, affecting approximately 20\% to 30\% of premenopausal women, accounting for over 20\% of all gynaecological consultations. Primary amenorrhoea is rare, occurring in less than 1\% of the population, whereas secondary amenorrhoea (excluding pregnancy) has an estimated prevalence of 3\% to 4\%. In Australia, gynaecological presentation rates align with international statistics, with heavy bleeding and endometriosis being major drivers of specialist referrals.

\section*{Aetiology and Pathophysiology}
The underlying causes of menstrual disorders are diverse and are clinically classified depending on the specific presentation.
\begin{itemize}
    \item \textbf{PALM-COEIN classification for abnormal uterine bleeding:} Developed by the International Federation of Gynecology and Obstetrics (FIGO). The structural causes (\textbf{PALM}) comprise endometrial \textbf{P}olyps, \textbf{A}denomyosis (endometrial tissue within the myometrium), \textbf{L}eiomyomas (uterine fibroids), and \textbf{M}alignancy or hyperplasia. The non-structural causes (\textbf{COEIN}) include \textbf{C}oagulopathy (e.g. von Willebrand disease), \textbf{O}vulatory dysfunction (e.g. polycystic ovary syndrome [PCOS]), \textbf{E}ndometrial dysfunction, \textbf{I}atrogenic causes (e.g. intrauterine devices, anticoagulants), and those \textbf{N}ot yet classified.
    \item \textbf{Prostaglandin-mediated dysmenorrhoea:} In primary dysmenorrhoea, the sloughing of the endometrium releases excess prostaglandins (specifically PGF$_{2\alpha}$ and PGE$_2$), which stimulate intense myometrial contractions, causing uterine vasoconstriction, ischaemia, and pain.
    \item \textbf{Secondary dysmenorrhoea pathology:} Pain caused by underlying pelvic pathology, most commonly endometriosis (ectopic endometrial implants triggering local inflammation and adhesions), adenomyosis, or pelvic inflammatory disease (PID).
    \item \textbf{Hypothalamic-pituitary-ovarian (HPO) axis dysfunction:} Amenorrhoea or oligomenorrhoea arising from disruption of pulsatile GnRH release. This can be secondary to functional hypothalamic amenorrhoea (driven by severe psychological stress, low body weight, or intensive athletic training), hyperprolactinaemia, or premature ovarian insufficiency (POI).
\end{itemize}

\section*{Clinical Features}
The clinical presentation of menstrual disorders depends on the specific cycle disturbance and the presence of pelvic pathology.
\begin{itemize}
    \item \textbf{Heavy menstrual bleeding:} Soaking through one or more sanitary pads or tampons every hour for several consecutive hours, the need for double sanitary protection, passing blood clots larger than a 50-cent coin, or bleeding lasting longer than 7 days.
    \item \textbf{Dysmenorrhoea:} Sharp, cramping lower abdominal pain starting just before or at the onset of menstruation, often radiating to the lower back and thighs, and potentially accompanied by nausea, diarrhoea, headache, and fatigue.
    \text{\item \textbf{Amenorrhoea:}} The complete absence of menarche by age 15 (primary amenorrhoea) or the cessation of established menses for 3 to 6 months or longer (secondary amenorrhoea).
    \item \textbf{Intermenstrual or post-coital bleeding:} Vaginal bleeding occurring between regular periods or immediately after sexual intercourse, which can indicate cervical or endometrial polyps, infection, or dysplasia.
    \item \textbf{Systemic symptoms of anaemia:} Fatigue, generalised weakness, pale skin (pallor), shortness of breath, and palpitations resulting from chronic iron deficiency secondary to menorrhagia.
    \item \textbf{Hyperandrogenic symptoms:} Hirsutism (excessive facial/body hair), severe acne, and weight gain, which are clinical signs of underlying ovulatory dysfunction such as PCOS.
\end{itemize}

\section*{Differential Diagnosis}
Clinicians must carefully differentiate between various menstrual, reproductive, and pelvic pathologies during evaluation.
\begin{itemize}
    \item \textbf{Primary vs.\ secondary dysmenorrhoea:} Primary dysmenorrhoea typically begins shortly after menarche, has no organic pathology, and responds well to NSAIDs. Secondary dysmenorrhoea onset is usually later in life, worsens progressively, and is associated with deep dyspareunia (painful sex) and chronic pelvic pain (suggesting endometriosis or adenomyosis).
    \item \textbf{Abnormal uterine bleeding vs.\ pregnancy complications:} Differentiating heavy menses from miscarriage, threatened abortion, or an ectopic pregnancy, which can present with irregular vaginal bleeding and pelvic pain.
    \item \textbf{Anovulatory bleeding vs.\ bleeding disorders:} Distinguishing the irregular, heavy bleeding of anovulatory cycles (common in adolescence and perimenopause) from a primary bleeding disorder (e.g. von Willebrand disease or platelet dysfunction).
    \item \textbf{Hypothalamic amenorrhoea vs.\ premature ovarian insufficiency:} Differentiating absent periods due to low body weight/stress (low FSH, low oestrogen) from early menopause (high FSH, low oestrogen).
\end{itemize}

\section*{When to Seek Medical Care}
Patients should seek medical evaluation for menstrual symptoms that interfere with their daily activities, severe pain not managed by over-the-counter pain relievers, irregular bleeding between periods, or if periods have stopped for more than three months.

\textbf{Call 000 (or local emergency services) for:}
\begin{itemize}
    \item Severe, sudden, unilateral lower abdominal or pelvic pain, which may indicate an ectopic pregnancy, ovarian torsion, or a ruptured ovarian cyst.
    \item Extremely heavy vaginal bleeding, such as soaking two or more pads per hour for consecutive hours, accompanied by dizziness, lightheadedness, confusion, or fainting (haemorrhagic shock).
    \item High fever, severe pelvic pain, and foul-smelling vaginal discharge, suggesting pelvic inflammatory disease or toxic shock syndrome.
\end{itemize}

\section*{Diagnosis and Investigations}
A structured diagnostic approach is essential to classify the menstrual disorder and guide targeted therapy.
\begin{itemize}
    \item \textbf{Pregnancy test:} A urinary or serum beta-hCG test is the mandatory first step in evaluating any reproductive-aged individual with menstrual changes or amenorrhoea.
    \item \textbf{Full blood count (FBC) and iron studies:} Assessing haemoglobin, haematocrit, ferritin, and transferrin saturation to diagnose and quantify iron-deficiency anaemia.
    \item \textbf{Transvaginal pelvic ultrasound (TVUS):} The primary imaging modality to assess uterine size, locate fibroids, measure endometrial thickness, and identify features of adenomyosis or polycystic ovaries.
    \item \textbf{Hormonal profile:} Serum tests for thyroid-stimulating hormone (TSH), prolactin, follicle-stimulating hormone (FSH), luteinising hormone (LH), estradiol, and free testosterone to evaluate HPO axis function.
    \item \textbf{Coagulation screen:} Testing prothrombin time (PT), activated partial thromboplastin time (aPTT), and von Willebrand factor antigen/activity, especially in adolescents presenting with menorrhagia at menarche.
    \item \textbf{Endometrial biopsy:} Indicated in women aged 45 and older, or younger women with risk factors (e.g. obesity, chronic anovulation), to rule out endometrial hyperplasia or endometrial carcinoma.
\end{itemize}

\section*{Management}
Management involves a combination of medical, surgical, and lifestyle interventions tailored to the patient's symptoms, underlying pathology, and reproductive goals.
\begin{itemize}
    \item \textbf{Hormonal therapies:} The levonorgestrel-releasing intrauterine system (LNG-IUS, e.g. Mirena) is the gold standard for reducing heavy bleeding. Combined oral contraceptive pills (COCP), progestogen-only pills, or implants are effective for regulating cycles and managing primary dysmenorrhoea and endometriosis.
    \item \textbf{Antifibrinolytics:} Oral tranexamic acid taken only during bleeding days to reduce menstrual blood loss by stabilizing clots, suitable for individuals not wanting hormonal therapy.
    \item \textbf{Non-steroidal anti-inflammatory drugs (NSAIDs):} Prostaglandin synthetase inhibitors, such as mefenamic acid, naproxen, or ibuprofen, taken at the onset of menstruation to relieve pain and reduce flow.
    \text{\item \textbf{Surgical interventions:}} Hysteroscopic polypectomy or myomectomy to remove structural lesions. Endometrial ablation or uterine artery embolisation for menorrhagia. Laparoscopic excision of endometrial implants for endometriosis. Hysterectomy is the definitive treatment for adenomyosis or fibroids when childbearing is complete.
    \item \textbf{Iron replacement:} Oral or intravenous iron supplementation to correct iron deficiency and anaemia.
    \item \textbf{Non-pharmacological measures:} Applying local heat (heat packs), acupuncture, and lifestyle modifications (dietary adjustments, regular exercise) to support symptom management.
\end{itemize}

\section*{Complications}
Untreated menstrual disorders can lead to long-term physical, reproductive, and metabolic complications.
\begin{itemize}
    \item \textbf{Severe iron-deficiency anaemia:} Chronic, severe depletion of red blood cells, leading to cardiovascular strain, cognitive impairment, and profound fatigue.
    \item \textbf{Infertility and subfertility:} Resulting from chronic anovulation (PCOS, POI), pelvic scarring and tubal obstruction (endometriosis, PID), or structural uterine cavities (fibroids, polyps).
    \item \textbf{Endometrial hyperplasia and carcinoma:} Caused by prolonged, chronic exposure of the endometrium to unopposed oestrogen in patients with long-standing anovulation (e.g. untreated PCOS).
    \item \textbf{Osteopenia and osteoporosis:} Premature bone density loss resulting from prolonged hypo-oestrogenism associated with functional hypothalamic amenorrhoea or POI.
    \item \textbf{Psychosocial impact:} High rates of anxiety, depression, and social withdrawal due to chronic pelvic pain or fear of public bleeding accidents.
\end{itemize}

\section*{Prognosis and Outcomes}
The prognosis for individuals with menstrual disorders is generally excellent, with a wide array of highly effective medical and surgical treatments available. The majority of patients with primary dysmenorrhoea and menorrhagia experience complete or significant symptom control with conservative medical therapies, such as the LNG-IUS or oral contraceptives. Primary dysmenorrhoea often improves naturally with age and following childbirth. Secondary disorders, such as endometriosis and adenomyosis, require long-term, multi-disciplinary management, but symptoms typically resolve completely following the menopause.

\section*{Prevention and Self-Care}
While many menstrual disorders cannot be directly prevented, self-care strategies help manage symptoms and promote general reproductive health.
\begin{itemize}
    \item \textbf{Menstrual cycle tracking:} Keeping a detailed record of cycle dates, bleeding duration, flow intensity, and pain levels using a calendar or mobile application.
    \item \textbf{Maintain a healthy body weight:} Avoiding extreme weight fluctuations, severe caloric restriction, or excessive physical exertion, which can disrupt the HPO axis.
    \item \textbf{Dietary iron optimization:} Consuming iron-rich foods (lean red meat, dark leafy greens, fortified cereals) along with vitamin C to enhance absorption.
    \item \textbf{Stress management:} Incorporating relaxation techniques, such as yoga, meditation, and cognitive behavioural therapy, to maintain normal hormonal cycles.
    \item \textbf{Heat therapy:} Using hot water bottles or adhesive heat patches on the lower abdomen to relax uterine muscles during painful cramps.
\end{itemize}

\section*{Sources and Further Reading}
The following organisations provide reliable medical information and support for menstrual health:
\begin{itemize}
    \item Royal Australian and New Zealand College of Obstetricians and Gynaecologists (RANZCOG). \textit{Patient education resources on heavy bleeding and pain.} https://ranzcog.edu.au/
    \item Jean Hailes for Women's Health. \textit{Comprehensive resources on periods, PCOS, and endometriosis.} https://www.jeanhailes.org.au/
    \item Healthdirect Australia. \textit{Overview of menstrual cycle issues and treatments.} https://www.healthdirect.gov.au/
    \item Mayo Clinic. \textit{Menstrual disorders: Diagnosis, causes, and therapies.} https://www.mayoclinic.org/
    \item NHS. \textit{Periods: Common problems and management.} https://www.nhs.uk/
\end{itemize}
